% !TEX root = ../thesis.tex
% =============================================================================
% Chapter 6: Discussion
% =============================================================================

\chapter{Discussion}
\label{ch:discussion}

This chapter interprets the four-stage results in the light of the research questions posed in \cref{ch:introduction}, identifies the limitations that bound the conclusions, and outlines directions for future work. \Cref{sec:discussion-findings} collects the main findings on whether trajectory-based behavioural feedback guides LLMs to design better metaheuristics. \Cref{sec:discussion-power} addresses the statistical power that the seed budget afforded the head-to-head comparison. \Cref{sec:discussion-bias} examines the CMA-ES bias exhibited by Gemini~3 Flash and the consequences of that bias for behavioural feedback specifically. \Cref{sec:discussion-channel} contrasts proximal and distal feedback channels and explains why prescriptive behavioural feedback steered the LLM in only a minority of cases. \Cref{sec:discussion-steering} examines why successful steering, when it occurred, often failed to translate into higher AOCC. \Cref{sec:discussion-complexity} interprets the AOCC-complexity correlation as partly an artefact of how the loop accumulates code over generations. \Cref{sec:discussion-novelty} considers why every Stage~4 winner is a CMA-ES variant and what would be needed to push the loop beyond canonical recombination.

% ===========================================================================
\section{Main Findings}
\label{sec:discussion-findings}

The headline result of Stage~4 is that feedback raises the floor of LLaMEA performance rather than its ceiling. Combined feedback's seed-to-seed standard deviation is roughly five times smaller than vanilla's, and a Levene test against vanilla returns $p = 0.030$, the only Stage~4 contrast to clear $\alpha = 0.05$. The mean-AOCC lift over vanilla is directionally consistent across all three feedback conditions but does not reach significance under Holm-corrected one-sided Welch tests at the same 10-seed budget. Two further patterns reinforce the variance interpretation. At the per-instance level, the three feedback conditions split 19 of 20 leaderboard wins between them while vanilla wins only one. At the per-seed level, vanilla produces four runs that finish below $0.88$ against two such runs across the 30 feedback-augmented seeds combined. Both describe a tighter performance distribution under feedback rather than a higher mean.

Together these signals suggest that additional feedback channels do not so much expand what a strong LLM can produce as suppress its worst failures. Vanilla AOCC-only feedback occasionally lets the LLaMEA loop settle into a poor-performing region of code space from which it does not recover within the 500-candidate budget. Behavioural and structural feedback both appear to help the loop avoid those attractors. The fact that combining the two channels does not produce additive gains over either alone, despite both individually outperforming vanilla, is consistent with their carrying overlapping rather than complementary information about what makes an algorithm good. Behavioural features describe the search dynamics that the code produces, structural features describe the code itself, and on this benchmark both appear to push the LLM toward a similar, more elaborate, CMA-like region of design space.

The overall assessment is therefore qualified but positive. Behavioural feedback does guide LLMs to design better optimisation algorithms, but the improvement manifests primarily as reduced failure-mode variance and reliable per-instance wins rather than as a large mean shift. The full-suite re-evaluation in \cref{sec:results-stage4-generalisation} reinforces this. All four LLM-generated winners outperform a CMA-ES baseline on the full $1{,}000$-instance MA-BBOB pool across $d \in \{5, 10, 20\}$, with \texttt{BIPOP\_MA\_CMA\_ES} from the neutral condition leading at $d = 10$ and $d = 20$ and remaining within $0.015$ AOCC of the leader at $d = 5$.

A notable limitation on these conclusions is the rule used to select the five behavioural features carried into Stage~4. We took the top five Stage~3 conditions by mean AOCC, which guaranteed that the chosen features had performed well under at least one feedback framing but inherited the calibration mismatch documented in \cref{sec:results-drift}. An alternative selection rule, such as the top-$k$ by Stage~3 Spearman correlation recomputed on Gemini~3 Flash output, or by a joint AOCC-and-Spearman rank, would test the sensitivity of Stage~4 to the feature-selection criterion. We did not run this ablation. The four-condition $\times$ ten-seed Stage~4 design already exhausted our Gemini API budget, and the comparison would have required at least one further full Stage~4 run. We flag this as the most defensible single-experiment extension of the present work.

% ===========================================================================
\section{Statistical Power Analysis}
\label{sec:discussion-power}

The single largest limitation of this thesis is the resolution at which the seed budget allowed condition-vs-condition differences to be detected. Stage~3 ran 5 seeds per condition across 29 conditions. Stage~4 ran 10 seeds per condition across 4 conditions.

With only 5 seeds per group, the Mann--Whitney $U$ test has $\binom{10}{5} = 252$ possible rank interleavings, so the smallest two-tailed $p$-value it can ever produce is $2/252 \approx 0.008$. That minimum requires complete separation between the two groups, meaning every value in one group exceeds every value in the other. The expected min and max of $n = 5$ samples from a normal distribution sit about $\pm 1.16\,\sigma$ from the mean\footnote{For $n$ i.i.d. samples from $\mathcal{N}(0, \sigma^2)$, the expected largest order statistic is $\sigma\, E[Z_{(n)}]$, where $Z_{(n)} = \max(Z_1, \dots, Z_n)$ for standard normal $Z_i$. With $E[Z_{(n)}] = \int_{-\infty}^{\infty} x \cdot n\,\Phi(x)^{n-1} \phi(x)\,\mathrm{d}x$ and no closed form, numerical integration at $n = 5$ gives $E[Z_{(5)}] \approx 1.163$. The expected smallest is $-1.163$ by symmetry.}. Complete separation thus requires the higher group's expected minimum ($\mu_A - 1.16\sigma_A$) to exceed the lower group's expected maximum ($\mu_B + 1.16\sigma_B$), equivalent to a mean gap of at least $1.16(\sigma_A + \sigma_B)$. For vanilla (seed-to-seed AOCC standard deviation $\sigma_v = 0.108$) versus the best-performing Stage~3 condition \texttt{neutral-intensification\_ratio} ($\sigma_b = 0.021$) that threshold is about $0.15$, well above the observed $+0.045$ gap. The test therefore cannot rule out chance even though the direction is consistent.

Doubling to 10 seeds per condition, as in the Stage~4 setup, raises the rank interleavings to $\binom{20}{10} = 184{,}756$. The expected min and max of $n = 10$ normal samples sit about $\pm 1.54\,\sigma$ from the mean (the $n = 10$ analogue of the $1.16$ value above), so the complete-separation threshold becomes $1.54(\sigma_A + \sigma_B)$. Even for the tightest Stage~4 pair, vanilla ($\sigma_v = 0.048$) versus combined ($\sigma_c = 0.009$) (\cref{tab:phase4-final-aocc}), this works out to a threshold of roughly $0.087$, and the comparisons against neutral and SAGE land closer to $0.13$. The observed Stage~4 mean lift of feedback over vanilla is only $\sim 0.02$ AOCC, still well below all three thresholds, which is why the Holm-corrected Welch tests of \cref{sec:results-stage4-agg} fall short of significance even though the direction is consistent across all three feedback conditions.

The seed budget itself was constrained almost entirely by Gemini API cost. The 29-condition feature-by-format sweep at 100 candidates per seed already amounted to $14{,}500$ LLaMEA iterations, and doubling to 10 seeds would have cost $29{,}000$ iterations, larger than the entire Stage~4 budget reserved for the head-to-head comparison. Stage~4 itself consumed $20{,}000$ candidate algorithms across four conditions and a sizeable Gemini API spend. Within the available budget the trade-off was either many feature-format conditions at low resolution (Stage~3) or fewer conditions at higher resolution (Stage~4), and we chose to spend our seeds where the design comparison was most critical.

A natural alternative is to step away from closed-API models entirely. Stage~1 showed that medium-sized open-weight models in the 24--32B parameter range, particularly Qwen3.5~27B and Devstral Small~2 (24B), can generate high-quality metaheuristics, although they trail the Gemini family in both AOCC and reliability. The next obvious step is to evaluate larger open-weight models in the 70--200B parameter range. Such models can be served on university hardware at marginal per-token cost, which would lift the seed budget by an order of magnitude or more and bring the kind of mean differences observed in Stage~4 within statistical reach. They would also permit the full $10 \times 5$ seed-by-condition factorial design that the Gemini budget could not support.

% ===========================================================================
\section{Model Bias}
\label{sec:discussion-bias}

A qualiative finding of Stage~4 is that Gemini~3 Flash exhibits a strong CMA-ES family bias regardless of the feedback string it receives. Every best-found algorithm across all four conditions, vanilla, neutral, SAGE, and combined, is a CMA-ES variant, and \cref{sec:results-stage4-winners} catalogues both the canonical mechanisms each winner inherits verbatim and the surface decorations layered on top. A reasonable interpretation is that Gemini~3 Flash, has a strong prior over which optimisation-algorithm code patterns are ``correct'', and that prior is anchored on the CMA-ES implementation. CMA-ES is a well-established strong algorithm on this class of benchmark, so the bias is empirically defensible. Converging on a CMA-ES variant is not a failure mode of the LLM but a constraint on what current feedback can induce it to produce.

One specific consequence of this bias is that it blunted Stage~3 comparative feedback by collapsing the reference values that format relies on. The Stage~1 dataset, drawn from ten heterogeneous models, populated a wide region of behavioural space and produced reference values for top-10\% and bottom-25\% performers that spanned several decades for some features. When Stage~3 deployed only Gemini~3 Flash, the behavioural distribution narrowed sharply. As \cref{sec:results-drift} documents, the bottom-25\% values for several features rose by an order of magnitude or more between stages while the top-10\% values barely moved, because the Stage~1 bottom was populated mostly by weaker open-weight models whereas the Stage~3 bottom is still produced by Gemini Flash. The gap between the two tiers, which determines the normalisation scale for comparative feedback, therefore collapsed by an order of magnitude for several features, and the distance signals the comparative format sends became small relative to the actual feature dispersion in the Stage~3 generation. The implication for future work is that behavioural-feedback calibration must be model-specific. A target-model calibration pass, screening the chosen LLM on its own to derive feature ranks and tier boundaries before the feedback experiment, would tighten the signal that prescriptive feedback attempts to convey.

% ===========================================================================
\section{Proximal versus Distal Feedback Channels}
\label{sec:discussion-channel}

\Cref{sec:results-decoupling} reported that prescriptive feedback moved the median feature value in the advised direction relative to neutral in only 7 of 19 cases, a steering accuracy of $37\%$. This is low enough to demand explanation, particularly because LLaMEA-SAGE~\cite{vanStein2026LLaMEASAGE} demonstrates that current-generation LLMs reliably comply with prescriptive natural-language advice when the target is something they can directly modify. SAGE issues prescriptive advice on structural features such as cyclomatic complexity and AST shape, and the resulting algorithms do shift on those axes with measurable AOCC gains. The contrast with our $37\%$ steering rate on behavioural features therefore points at a structural difference rather than a model limitation.

The relevant difference is the causal channel. SAGE's structural advice is proximal. The LLM edits the abstract syntax tree, and the next AST is a function of those edits, so compliance is mechanical. Behavioural advice is distal. The LLM still edits code, but the behavioural feature emerges from running that code on a benchmark, and the mapping from a code change to a behavioural delta is uncertain even to a human author. An LLM that genuinely tries to comply with prescriptive behavioural advice must therefore guess at this mapping, and on a non-trivial mutation that guess is unreliable. The $37\%$ rate is what we should expect from a steering channel that is causally indirect.

This framing has a concrete implication for future feedback designs. Behavioural features are not in themselves unsteerable, but they cannot be steered by direct prescription the way structural features can. They must instead be steered through their structural antecedents. A prescription like ``increase the intensification ratio'' asks the LLM to control a measurement it cannot directly write. Guidance like ``narrow the sampling distribution around the incumbent during the second half of the budget'' asks it to control structural code that, if executed, would produce the desired behavioural shift. Translating behavioural targets into structural prescriptions before issuing them is the natural next step, and would test whether the steering deficit observed here is a property of behavioural features per se or only of how they were prescribed.

% ===========================================================================
\section{Why Successful Steering Did Not Translate to AOCC}
\label{sec:discussion-steering}

Even when prescriptive feedback successfully steered a behavioural feature, it often failed to lift AOCC. \texttt{x\_spread\_early} and \texttt{fitness\_autocorrelation} are both steered correctly by both prescriptive formats yet achieve lower AOCC than neutral. The previous section explains why prescriptive formats often fail to act on the LLM. It does not explain why, when they do act, the result is no better and sometimes worse. Two explanations are plausible, and we suspect both contribute.

The first concerns the multivariate structure of the high-performance region. Good algorithms appear to be characterised by a combination of feature values rather than by any single feature in isolation. Pushing one feature toward its top-tier reference while leaving the others unconstrained can place a candidate in a region of behavioural space that no actual top-tier algorithm occupies, so single-feature optimisation can take the LLM off the high-performance manifold rather than along it. Future feedback designs should target behavioural profiles rather than individual features. One concrete instantiation would be to identify clusters of high-performing algorithms in behavioural space (the upper-left quadrant of \cref{fig:tsne-behavioral} is a candidate), report the candidate's nearest-cluster centroid as the target, and frame the advice as ``move toward this configuration of features'' rather than ``increase this one feature.''

The second concerns the granularity at which behavioural features are computed. Each feature reported in the feedback prompt is averaged across 50 MA-BBOB runs (10 functions $\times$ 5 instances) in Stage~3 and 100 runs in Stage~4. This compresses behaviour across landscapes that differ substantially in conditioning, modality, and separability into a single scalar. An algorithm may need a high intensification ratio on unimodal high-conditioning functions and a low one on multimodal weak-structure functions, and reporting the cross-instance average obscures that conditional structure. A natural refinement, again left to future work, is to report behavioural features per MA BBOB function and let the LLM see that its algorithm exploits well on separable functions but explores too aggressively on multimodal ones. Furthermore, behavioural feedback may need to be studied at the instance or group level before it can be aggregated effectively.

A third candidate cause, which we examined and reject, is calibration mismatch. The prescriptive reference values were derived from the heterogeneous Stage~1 model pool but deployed against Gemini~3 Flash alone, whose behavioural distribution is much narrower. The natural worry is that even successful steering pushes a candidate to a reference value that sits at the wrong scale for this deployment model, so the shift overshoots or lands in a region where the original AOCC-feature relationship no longer holds. Inspection of the recomputed Stage~3 tier values, however, shows that for most features the prescriptive targets sit close to where Gemini's own distribution places its top-tier candidates, so successful steering does not place candidates in a fundamentally different region of behavioural space. Calibration mismatch is a real issue of the prescriptive design and plausibly degrades the ability of comparative feedback to correctly steer an algorithm to a specific behavioural value. Once steering succeeds, the mismatch ceases to act. The candidate has been pulled to a reference value that sits close to where Gemini's own top-tier algorithms already sit, so the original calibration error has no remaining channel through which to depress AOCC. We initially suspected this explanation but reject it on closer examination.

The fact that neutral framing, which sidesteps the two contributing problems by simply reporting values without imposing an objective, consistently outperforms the prescriptive variants is consistent with this view. Neutral feedback adds informative context to the prompt without asking the LLM to optimise something it should not be optimising in isolation. The takeaway is that using the methodology of this thesis, describing the candidate's behaviour to the LLM produces better algorithms than instructing it how to change that behaviour, but we should not entirely throw away using target behavioural profiles as a way to guide an LLM to better performance.

% ===========================================================================
\section{Code Complexity}
\label{sec:discussion-complexity}

\Cref{sec:results-ast} showed that neutral feedback produces structurally more complex code than directional or comparative feedback, and that across the full Stage~3 population code complexity correlates positively with AOCC. Total AST nodes ($\rho = +0.48$), NumPy calls ($+0.60$), and numeric constants ($+0.66$) are among the strongest individual predictors of algorithm quality. This reproduces the central finding of LLaMEA-SAGE~\cite{vanStein2026LLaMEASAGE}, that complexity-promoting structural feedback improves performance, and offers a candidate explanation for why neutral behavioural feedback works. It shifts the LLM toward writing more elaborate code on axes that empirically correlate with stronger optimisation behaviour.

The interpretation of this complexity deserves scrutiny. Two observations from Stage~4 suggest that the AOCC-complexity correlation is partly an artefact of how the LLaMEA loop accumulates code rather than a sign that complexity itself drives quality. The first is the catalogue of dead code in the Stage~4 winners reported in \cref{sec:results-stage4-winners}. The second is the failure-rate climb in \cref{fig:phase4-failure-bygen}. SAGE's failure rate rises from $12.8\%$ in candidates 100--199 to $27.9\%$ in candidates 400--499, and combined climbs from $16.8\%$ to $23.7\%$ over the same window, even though the prompt template does not change across generations.

A single mechanism plausibly produces both. LLaMEA's mutation prompt instructs the LLM to ``refine the strategy of the selected solution to improve it'', with no mention of simplification, pruning, or trimming. Conditional on the LLM treating ``refine'' as an instruction to add rather than to rewrite, the body of code under mutation grows monotonically across generations. The dead-code accumulation is the visible content of that growth. The failure-rate climb is its statistical consequence. Each marginal mutation lands on a larger surface where any single error fails the whole candidate, suggesting that LLMs struggle to one-shot working improvements to complex code. Under this reading, the AOCC-complexity correlation does not show that more elaborate code is better. It shows that the loop produces more elaborate code over time, and that AOCC also rises over time, with the underlying cause being whatever the LLM is doing well rather than the elaboration itself.

Two future-work directions address how the loop accumulates code. The first is diagnostic. Whether each marginal LLaMEA mutation adds executed code or only decorative tokens can be answered by AST-diffing consecutive candidates and tracing which added subtrees are actually executed during evaluation. This would discriminate between meaningful and inert complexity at the per-mutation level and turn the dead-code observations into a quantitative signal. The second is corrective. Mutation acceptance could be gated on a measured behavioural delta from the parent rather than on AOCC alone, so that mutations which only relabel or restructure code without changing what it does are filtered before they accumulate. The mutation prompt itself could also be modified to explicitly reward trimming as well as additions, replacing ``refine'' with language that admits both directions of edit.

% ===========================================================================
\section{Algorithmic Novelty}
\label{sec:discussion-novelty}

The novelty of the generated metaheuristics is also in question. Every Stage~4 winner is a CMA-ES variant with two or three surface decorations layered over a textbook implementation, and the decorations themselves are standard variants from the CMA-ES literature: mirrored sampling~\cite{mirrored2010}, active CMA with negative rank-$\mu$ weights~\cite{active2006}, and IPOP and BIPOP restart schemes~\cite{ipop2005, bipop2009}. The momentum-style mean-shift accumulation in \texttt{OM\_CMA\_EIS} is not a named published variant but is conceptually close to the cumulation already present in the canonical evolution path $p_c$. The LLM is recombining well-known mechanisms rather than discovering new ones, and the LLaMEA loop in its current form functions more as a sophisticated hyperparameter search across known mechanisms than as a design-space explorer in the sense Hoos~\cite{hoos2012pbo} originally envisaged.

Some of this concentration reflects how hard genuine algorithmic novelty is. Decades of expert research have produced relatively few mechanisms that materially advance the state of the art on continuous black-box optimisation and many published metaheuristics rediscover algorithmic components under new names~\cite{aranha2022metaphor}. The bar for novelty in this domain is not ``produce code that has not been written before'' but ``produce a mechanism that improves on a strong, mature baseline'', and that bar is high enough that current LLMs may not be able to clear it regardless of how the loop is configured. The recombination behaviour observed here is consistent with what one would expect from a system that knows the canonical mechanisms well and lacks the capacity to invent new ones.

With that said, two structural features of the LLaMEA loop may reinforce a model's bias, though we did not run the ablations that would be needed to establish either as a cause, and we offer them as conjectures rather than findings. The first is the selection rule. The $(1{+}1)$ elitist replacement keeps an offspring only if it matches or beats the parent on AOCC, which could in principle penalise a novel design that is not already competitive with the incumbent on the generation it appears, since the loop offers no protected niche in which such a design might mature. This intuition should be read with care. Unlike a conventional hillclimber, the LLaMEA mutation operator is not confined to local moves. Under the 10\% explore prompt the LLM can propose a wholesale new design rather than a refinement of the incumbent. So the loop does have a channel through which a genuinely different algorithm can enter, and whether elitist selection nonetheless suppresses novelty on net is an empirical question we leave open. The second is the in-context example. Each mutation prompt shows the LLM the current best algorithm as the example to work from, which after the first few generations is already a CMA-ES variant. Whether this anchors subsequent generations to the CMA-ES family, or merely co-occurs with a prior the LLM would hold regardless, likewise cannot be settled without an ablation that varies the example independently of the incumbent.

We also note that the $(1{+}1)$ configuration studied here is not the only option. More recent LLaMEA variants adopt larger populations with comma-selection $(\mu, \lambda)$ strategies and explicit diversity-preservation mechanisms~\cite{vanStein2025Behaviour}, which relax precisely the selection pressure conjectured above and would be the natural setting in which to test whether it suppresses novelty.

We identify two complementary directions could push the loop beyond canonical recombination. The first is to expand the LLM's effective context with a corpus of past LLaMEA-generated algorithms via retrieval-augmented generation~\cite{rag2020,ragcode2021}. The $16{,}538$ valid Stage~4 algorithm-feature pairs produced by this thesis, together with the tens of thousands more available on public LLaMEA Zenodo repositories, form a candidate corpus. Each mutation prompt would include a handful of high-AOCC exemplars deliberately drawn from non-CMA regions of design space, biasing the LLM away from its default basin without forbidding it outright. The second is more direct. The mutation prompt could forbid canonical mechanisms by name and require the LLM to justify novelty against a list of known variants. This trades the implicit nudge of retrieval-augmented exemplars for an explicit constraint, at the cost of more brittle prompting.
